\section{Plan de elicitación}\label{sec:elicitacion}

\subsection{Criterios de selección de técnicas}
La entrevista recoge explicaciones que el responsable puede expresar. La observación ayudará
a reconocer pasos que realiza de forma habitual y quizá no menciona. El cuestionario permitirá
recoger la experiencia de clientes con preguntas breves. Se combinan fuentes y técnicas según
el conocimiento disponible~\cite{pohl2025,ireb2026}.

La Tabla~\ref{tab:tecnicas} distingue las fuentes consultadas de las actividades propuestas.
Solo se ha realizado la entrevista SE-01; las demás actividades requieren coordinación y el
consentimiento correspondiente.

\begingroup\small\setstretch{1.1}
\begin{longtable}{@{}L{1cm}L{2.7cm}L{2.4cm}L{5.7cm}L{1.8cm}@{}}
\caption{Técnicas y fuentes por rol.}\label{tab:tecnicas}\\
\toprule
\textbf{Rol} & \textbf{Técnica} & \textbf{Conoci\-miento} & \textbf{Motivo y estado} & \textbf{Instru\-mento} \\
\midrule\endfirsthead
\toprule
\textbf{Rol} & \textbf{Técnica} & \textbf{Conoci\-miento} & \textbf{Motivo y estado} & \textbf{Instru\-mento} \\
\midrule\endhead
\bottomrule\endfoot
PI-01 & Entrevista & Explícito: tareas y prioridades. & SE-01 realizada; permite explicar decisiones y casos concretos. & ENT-01 \\
PI-02 & Observación & Habitual o tácito: pasos de recepción y entrega. & Prevista para contrastar qué se anota y cómo se localiza el servicio. & OBS-01 \\
PI-03 & Observación & Habitual: actualización del estado. & Prevista; se observará el registro, sin intervenir equipos ni acceder a archivos privados. & OBS-01 \\
PI-04 & Revisión de ejemplos & Documentado: estado del cobro. & Prevista con registros anonimizados; solo se distinguirá pagado o pendiente. & OBS-01, O5 \\
PI-05 & Cuestionario & Explícito: experiencia del servicio. & Previsto; contrasta claridad de recepción, estado y entrega. No mide la línea base interna. & CUE-01 \\
PI-06 & Fuente indirecta & Declarado por el taller. & Se usa A4 y B5 de SE-01; no se consulta directamente mientras el módulo de inventario siga excluido. & ENT-01 \\
PI-07 & Revisión documental & Analítico: relaciones y decisiones. & El equipo contrasta entrevista, objetivos, modelos e instrumentos antes de integrar cambios. & Informe y EV-04 \\
PI-08 & Revisión de la guía & Normativo académico. & La guía define los criterios; las dudas se llevarán a tutoría, sin suponer aprobación previa. & Guía de E1 \\
\end{longtable}
\endgroup

\subsection{Cronograma de sesiones}
La Tabla~\ref{tab:sesiones} propone las siguientes actividades. Las semanas futuras son una
planificación del equipo; sus fechas dependen de la disponibilidad de quienes participen.

\begin{table}[H]\centering\small
\caption{Sesiones realizadas y previstas.}\label{tab:sesiones}
\noindent\begin{tabularx}{\linewidth}{@{}L{1cm}L{2.9cm}L{2.7cm}YL{2.5cm}@{}}
\toprule
\textbf{ID} & \textbf{Estado y momento} & \textbf{Rol / técnica} & \textbf{Propósito} & \textbf{Responsable} \\
\midrule
SE-01 & Realizada: 21/11/2026, 14:20; 12 min 44 s. & PI-01 / entrevista ENT-01. & Conocer proceso, dificultades y prioridades. & Jordan; notas de Jonathan. \\
SE-02 & Prevista: semana 5; fecha por acordar. & PI-02 a PI-04 / observación y ejemplos. & Contrastar registros, estados y datos para ON-01 a ON-03. & Jordan coordina; Jonathan registra. \\
SE-03 & Prevista: semana 5; fecha y participantes por acordar. & PI-05 / cuestionario. & Recoger experiencia del cliente sobre información y entrega. & Jordan coordina; equipo analiza. \\
SE-04 & Prevista: semana 6; fecha por acordar. & PI-01 / revisión conjunta. & Confirmar objetivos, alcance, exclusiones y correcciones. & Jonathan prepara; Jordan coordina. \\
\bottomrule
\end{tabularx}
\end{table}

SE-01 está resumida en el Anexo~\ref{anx:registro-entrevista}. Las nuevas sesiones requieren
explicar la actividad y sus permisos antes de comenzar. No se registran como resultados ya obtenidos.

\subsection{Instrumentos y trazabilidad}
Los anexos~\ref{anx:entrevista}, \ref{anx:cuestionario} y~\ref{anx:observacion} contienen los tres
instrumentos. La Tabla~\ref{tab:traza-instrumentos} muestra su relación con la propuesta actual.
La relación del guion con los objetivos se estableció después de analizar la entrevista; no
supone que esos objetivos ya estuvieran acordados al aplicarla.

\begin{table}[H]\centering\small
\caption{Relación de preguntas y observaciones con objetivos y alcance.}\label{tab:traza-instrumentos}
\noindent\begin{tabularx}{\linewidth}{@{}L{3.7cm}L{2.6cm}YL{2.4cm}@{}}
\toprule
\textbf{Preguntas / aspectos} & \textbf{Objetivo} & \textbf{Capacidad} & \textbf{Estado} \\
\midrule
ENT-01: B2, B4, C3--C4, D1 & ON-01, ON-02 & AL-01 a AL-03 & Aplicado. \\
ENT-01: B3, B7, C2, D1 & ON-03 & AL-04 & Aplicado. \\
ENT-01: E2 & ON-01 & AL-05 & Aplicado; alcance por confirmar. \\
CUE-01: C1--C4 & ON-02 & AL-01 a AL-03; percepción del cliente. & Sin aplicar. \\
CUE-01: C5 & ON-03 & AL-04; claridad del estado de pago. & Sin aplicar. \\
OBS-01: O1--O4 & ON-01, ON-02 & AL-01 a AL-03 & Sin aplicar. \\
OBS-01: O5 & ON-03 & AL-04 & Sin aplicar. \\
OBS-01: O6 & ON-01 & AL-05 & Sin aplicar. \\
\bottomrule
\end{tabularx}
\end{table}
